\section{Introduction} \label{intro}

Algorithms shape a growing share of consequential decisions. Many of these decisions carry consequences that fall, at least in part, on others. Physicians rely on diagnostic systems, firms let screening software rank job applicants, and financial advisors draw on algorithmic tools. In each setting, a human decides whether to rely on the algorithm. Who bears responsibility when the resulting decision causes harm is the subject of a lively legal \citep{lemley_remedies_2019, buiten_law_2023}, ethical \citep{matthias_responsibility_2004, santoni_de_sio_four_2021}, and policy \citep{european_commission_liability_2019,european_union_ai_act_2024} debate. Yet the normative question of where responsibility should lie is distinct from the behavioral questions of how responsibility is attributed and how those attributions shape behavior. Does delegating a decision to an algorithm change the responsibility that affected parties assign to the delegator? And does the prospect of being held responsible affect the decision to delegate? We provide experimental evidence on both questions.

For delegation between humans, the experimental literature points in a clear direction. Delegation shifts responsibility. Principals who delegate an unpopular decision to a human intermediary are punished less for it, feel less responsible for it, and delegate strategically for that reason \citep{bartling_shifting_2012, hamman_self-interest_2010, coffman_intermediation_2011, oexl_shifting_2013}. Recent evidence suggests that algorithms can inherit the intermediary's role. Delegating to an algorithm helps decision-makers escape blame for adverse outcomes \citep{feier_hiding_2022}, lets them shed moral responsibility for self-serving choices \citep{tontrup_strategic_2025} and makes dishonesty easier by obscuring intent \citep{kobis_delegation_2025}. If algorithms provide such cover, decision-makers who face potential punishment should be drawn to them, and delegation should protect the delegator when outcomes turn out badly. Yet no existing study varies whether delegators can be punished, leaving untested whether the prospect of punishment in fact motivates delegation.

We address both questions in a preregistered online experiment with 322 participants.\footnote{The preregistration (AsPredicted \#133980) is available at \url{https://aspredicted.org/hs9az8.pdf} and reproduced in full in Appendix~\ref{appendix:prereg}.} A decision-maker, Player~A, first completes ten rounds of a prediction task for her own payoff, and then makes an additional prediction that determines the payoff of a matched recipient, Player~B. Player~A can either make the consequential prediction for Player~B herself or delegate it to an algorithm calibrated to perform equally well at the task. Thus, the environment we study, a prediction task with outcome uncertainty performed by a decision-maker with no material stake in the affected party's outcome, resembles the settings in which algorithmic delegation typically operates in practice. We then vary between subjects whether Player~B can reduce Player~A's earnings. In the \textit{Punishment} condition, we elicit Player~B's punishment for every combination of delegation choice and prediction outcome, providing a within-subject test of whether delegation insulates Player~A from punishment, conditional on the outcome. In the \textit{No-Punishment} condition, no such punishment is possible. With everything except the possibility of punishment held fixed, any difference in delegation rates is attributable to that possibility, allowing us to test directly whether the prospect of being held responsible moves the delegation decision.

We find that the prospect of punishment reverses the prediction based on the literature. When punishment is possible, $42.5$ percent of decision-makers delegate to the algorithm. When it is not, $60.8$ percent do. The possibility of punishment thus \textit{reduces} delegation by $18.3$~percentage points. Delegation, in turn, offers no protection against punishment. Conditional on the outcome, Player~B punishes delegated and Player~A's own decisions similarly, and punishment responds almost entirely to the outcome of the prediction rather than to how it was made.

%old
%Thus, the possibility of punishment moves the delegation decision even though delegating draws no additional punishment. We explore several candidate explanations to reconcile these findings. A first candidate is anticipated punishment. Player~A may delegate less because she expects delegated decisions to be punished more harshly. Her elicited beliefs speak against this story. On average, her expected punishment is no higher for delegated than for self-made decisions, matching how Player~B actually punishes. A second candidate is punishment-induced effort. Facing potential punishment, Player~A may keep the decision in the hope of outperforming the algorithm. If so, non-delegators should improve on their earlier performance precisely when punishment is possible. Instead, improvement concentrates where punishment is impossible. What the data do show is selection. When punishment is possible, it is disproportionately the better-performing Player~As who keep the decision for themselves, and the delegation gap between conditions widens with performance. Both a preference for producing the good outcome oneself and a distaste for delegating uncertain decisions are consistent with this selection, but they operate in both conditions alike and cannot on their own generate the treatment difference. Lastly, not delegating can carry signaling value. By keeping the decision, Player~A may aim to convey to Player~B that she intends to try especially hard on his behalf, whereas delegating could be perceived as shirking the decision. Because the task is inherently uncertain, sending this signal requires no measurable extra effort. However, the signal carries monetary value only if Player~A expects apparent shirking to draw harsher punishment, which her beliefs do not show. Detached from actual punishment, she may instead care about Player~B's judgment itself, although this too explains the treatment difference only if the prospect of punishment is what activates this motive. Thus, while the drop in delegation under punishment is robust, which motive produces it is a question we narrow rather than settle.

The possibility of punishment thus affects the decision to delegate, even though delegating draws no differential punishment. We examine several explanations to reconcile these findings. The first is anticipated punishment. Player~A may keep the decision because she expects delegation to be punished more harshly. Her elicited beliefs provide no evidence of such an expectation and instead match how Player~B actually punishes. The second is punishment-induced effort. Player~A may keep the decision in hopes of outperforming the algorithm in the prediction task for Player~B, in which case decision-makers should improve on their earlier performance precisely when punishment is possible. Instead, improvement concentrates where punishment is impossible. Third, by keeping the decision, Player~A may aim to signal a willingness to try on Player~B’s behalf, whereas delegation may appear as shirking. Because performance is noisy, the choice to retain control can carry this signal even without a measurable improvement in performance. However, the signal carries monetary value only if Player~A expects apparent shirking to draw harsher punishment, which her beliefs do not indicate. Detached from actual punishment, Player~A may instead care directly about Player~B’s judgment. Because Player~B can form such a judgment in both conditions, however, this concern can explain the treatment difference only if the prospect of punishment activates or strengthens it. Lastly, the data show selection in performance. It is disproportionately the better-performing Player~As who retain the decision. Both an intrinsic preference for producing the good outcome oneself and a distaste for delegating uncertain decisions are consistent with this selection, but they operate in both conditions alike and cannot on their own generate the treatment difference. Thus, while the drop in delegation under punishment is robust, which motive produces it is a question we narrow down rather than settle.

%The possibility of punishment thus moves the delegation decision even though delegating draws no additional punishment. We examine several explanations to reconcile these findings. The first is anticipated punishment. Player~A may keep the decision because she expects delegation to be punished more harshly. Her elicited beliefs show no such expectation, and instead match how Player~B actually punishes. The second is punishment-induced effort. Player~A may keep the decision in the hope of outperforming the algorithm in the prediction task for Player~B, in which case decision-makers should improve on their earlier performance precisely when punishment is possible. Instead, improvement concentrates where punishment is impossible. What the data do show is selection. It is disproportionately the better-performing Player~As who keep the decision, and the delegation gap between conditions slightly widens with performance. Both an intrinsic preference for producing the good outcome oneself and a distaste for delegating uncertain decisions are consistent with this selection, but they operate in both conditions alike and cannot on their own generate the treatment difference. Lastly, by keeping the decision, Player~A may aim to merely convey to Player~B that she intends to try especially hard on his behalf, whereas delegating could be perceived as shirking the decision. Because the task is inherently uncertain, sending this signal requires no measurable additional effort. However, the signal carries monetary value only if Player~A expects apparent shirking to draw harsher punishment, which her beliefs do not show. Detached from actual punishment, she may instead care about Player~B's judgment itself, although this too explains the treatment difference only if the prospect of punishment is what activates this motive. Thus, while the drop in delegation under punishment is robust, which motive produces it is a question we narrow rather than settle.

Our contribution is twofold. First, we provide, to our knowledge, the first direct test of whether the prospect of punishment encourages delegation to an algorithm, as prior evidence on delegation to human intermediaries would predict \citep{bartling_shifting_2012, hamman_self-interest_2010, coffman_intermediation_2011, oexl_shifting_2013}. Our environment departs from these canonical designs in three respects. The intermediary is an algorithm rather than a person. The decision involves genuine uncertainty, so a bad outcome may reflect bad luck rather than bad intent. And its immediate consequences fall entirely on the recipient, so the underlying decision involves no tradeoff between the players’ payoffs and there is no self-serving choice to police. In this environment, the effect runs in the opposite direction. The prospect of punishment discourages rather than encourages delegation, and it does so selectively among the decision-makers most likely to succeed. This reversal is the paper's main empirical result, adding a determinant of preferences over algorithm use that operates through the decision context rather than through algorithm characteristics or the task.

Second, we use Player~B's punsihment decision to examine how an affected party assigns responsibility for algorithmic delegation when the delegation choice is uninformative about relative ability. The closest precedent is \citet{feier_hiding_2022}, whose treatments vary whether the intermediary is an algorithm or another participant under a fixed punishment regime, whereas we hold the intermediary fixed as an algorithm and vary whether punishment is possible. In their environment, decisions delegated to an algorithm are rewarded more generously than non-delegated decisions after bad outcomes. However, the same subject both makes a delegation choice and evaluates the decision, so that evaluations may partly reflect participants’ own delegation choices. Moreover, the delegation decision is strongly determined by beliefs about relative ability between oneself and the intermediary, inviting punishment of overconfidence rather than of delegation itself. Our design is intended to remove both channels by separating the delegator and evaluator roles and making delegation uninformative about relative ability. Once both channels are closed, the differential disappears, and punishment tracks the realized outcome instead.

The remainder of the paper is organized as follows. Section~\ref{literature} relates the study to the literature. Section~\ref{design} presents the experimental design and hypotheses. Section~\ref{results} reports the results and examines candidate mechanisms. Section~\ref{conclusion} concludes.

%Three design features distinguish our setting from the closest precedent, \citet{feier_hiding_2022}, two of them central to identification. First, the algorithm is calibrated at the \textit{individual} level rather than the session or group level. Existing designs typically expose subjects to a distribution of algorithmic performance and let them infer their own ability relative to that distribution; the delegation decision then carries strong signal value about the principal's self-perception, and any differential punishment can reflect overconfidence rather than responsibility shifting per se. By neutralizing relative-performance signaling at the individual level, we isolate the responsibility-shifting motive. Second, our \textit{No-Punishment} condition turns off the punishment possibility entirely, allowing us to identify the effect of anticipated punishment on the delegation decision itself---a test that none of the existing algorithmic-delegation designs permits. Third, we separate the delegator and evaluator roles across disjoint subject pools, rather than having the same subjects play both, removing the self-justification motives that a dual role can introduce into the punishment decision. Together these features sharpen the test of responsibility shifting in algorithmic delegation that the literature has thus far been able to deliver.

%These findings contribute to three strands of literature. First, we add to the experimental literature on responsibility avoidance through delegation \citep{bartling_shifting_2012,coffman_intermediation_2011,oexl_shifting_2013,hamman_self-interest_2010,hill_does_2015,steffel_passing_2016}, by extending the canonical setting to delegation to an algorithmic intermediary in a setting with genuine outcome uncertainty. Second, we contribute to the rapidly growing literature on algorithm aversion and adoption \citep{dietvorst_algorithm_2015,logg_algorithm_2019,burton_systematic_2020,mahmud_what_2022,chugunova_we_2022} by isolating an under-studied determinant of algorithm use: the institutional context in which the human user is held accountable. Third, and most directly, we build on \citet{feier_hiding_2022} and \citet{kirchkamp_sharing_2019} by providing a clean experimental test of two long-standing claims in the responsibility-and-algorithms literature, with a design that controls relative-performance signaling at the individual level and that includes a condition in which punishment is ruled out by construction. The contrast with \citet{feier_hiding_2022} is consequential: that study finds that decision-makers are held \textit{less} responsible for algorithmically delegated decisions, while we find that the punishment possibility \textit{reduces} delegation rather than increasing it. We argue that the two findings are reconcilable: the relative-performance signaling that drives much of Feier et al.'s differential punishment is closed off by construction in our design, isolating a different and milder margin (the motivational salience of accountability for the principal). Beyond the experimental literature, our findings speak to the regulatory debate over the assignment of responsibility for algorithmic decisions \citep{european_union_ai_act_2024}: within the bounds of the setting we study, holding human users accountable for algorithm-mediated outcomes appears to be a behavioral lever for direct human engagement with the decision, rather than merely a tool for ex-post moral attribution. We are explicit in Section~\ref{conclusion} about why the present evidence---a single online experiment on a stylised prediction task with one UK Prolific sample---supports this claim only in a bounded sense, and what would be required to extend it to the high-stakes algorithmic domains (medical, judicial, hiring, financial) that motivate the regulatory debate.
